% =============================================================
%  PREAMBLE — Common Lisp: Fondamenti e Pratica  (testo teorico)
%  Compilare con: pdflatex main  (x2)
% =============================================================

% --- Encoding e lingua ---
\usepackage[T1]{fontenc}
\usepackage[utf8]{inputenc}
\usepackage[italian]{babel}

% Blindatura caratteri accentati italiani per pdflatex+T1
% (evita problemi in contesti speciali come titoli tcolorbox)
\DeclareUnicodeCharacter{00E0}{\`a}  % à
\DeclareUnicodeCharacter{00E8}{\`e}  % è
\DeclareUnicodeCharacter{00E9}{\'e}  % é
\DeclareUnicodeCharacter{00EC}{\`\i} % ì
\DeclareUnicodeCharacter{00F2}{\`o}  % ò
\DeclareUnicodeCharacter{00F9}{\`u}  % ù
\DeclareUnicodeCharacter{00C0}{\`A}  % À
\DeclareUnicodeCharacter{00C8}{\`E}  % È
\DeclareUnicodeCharacter{00C9}{\'E}  % É
\DeclareUnicodeCharacter{00CC}{\`I}  % Ì
\DeclareUnicodeCharacter{00D2}{\`O}  % Ò
\DeclareUnicodeCharacter{00D9}{\`U}  % Ù
\DeclareUnicodeCharacter{00F3}{\'o}  % ó  (es. però)
\DeclareUnicodeCharacter{00FA}{\'u}  % ú

% --- Font ---
\usepackage{lmodern}
\usepackage{microtype}
\usepackage{graphicx}
\usepackage[outline]{contour}
\contourlength{0.6pt}

% --- Layout pagina: margini simmetrici, generosi ---
\usepackage[
  a4paper,
  top=2.8cm, bottom=3cm,
  left=3cm, right=3cm
]{geometry}

% --- Colori ---
\usepackage[dvipsnames,svgnames,table]{xcolor}

% Palette brand (condivisa col workbook)
\definecolor{clBg}{RGB}{248,248,248}
\definecolor{clKeyword}{RGB}{0,0,180}
\definecolor{clComment}{RGB}{110,110,110}
\definecolor{clString}{RGB}{0,120,0}

% Colori semantici
\definecolor{csBlue}{RGB}{30,80,160}
\definecolor{csGreen}{RGB}{34,110,34}
\definecolor{csOrange}{RGB}{200,100,0}
\definecolor{csPurple}{RGB}{130,50,150}
\definecolor{csTeal}{RGB}{0,110,110}
\definecolor{csRed}{RGB}{180,30,30}

% Box frame/back
\definecolor{boxConFrame}{RGB}{30,80,160}
\definecolor{boxConBack}{RGB}{235,241,252}
\definecolor{boxEsFrame}{RGB}{34,110,34}
\definecolor{boxEsBack}{RGB}{236,248,236}
\definecolor{boxWarnFrame}{RGB}{200,100,0}
\definecolor{boxWarnBack}{RGB}{255,249,220}
\definecolor{boxWbFrame}{RGB}{130,50,150}
\definecolor{boxWbBack}{RGB}{248,236,255}
\definecolor{boxKeyFrame}{RGB}{0,110,110}
\definecolor{boxKeyBack}{RGB}{230,248,248}
\definecolor{boxNoteFrame}{RGB}{130,130,130}
\definecolor{boxNoteBack}{RGB}{245,245,245}

% --- Matematica ---
\usepackage{amsmath}
\usepackage{amssymb}

% --- Listings ---
\usepackage{listings}

\lstdefinelanguage{CommonLisp}{%
  sensitive=false,
  morecomment=[l]{;},
  morestring=[b]",
  morekeywords=[1]{%
    defun, defmacro, defparameter, defvar, defconstant,
    defstruct, defclass, defmethod, defgeneric,
    let, let*, flet, labels, setf, setq,
    if, cond, when, unless, case, ecase,
    and, or, not,
    progn, prog1, prog2,
    loop, dotimes, dolist, do, do*,
    lambda, function,
    funcall, apply,
    mapcar, maplist, mapcan, mapc,
    remove-if, remove-if-not, find-if, find-if-not,
    every, some, notany, notevery,
    reduce, count-if, count,
    sort, stable-sort,
    return, return-from,
    values, multiple-value-bind, multiple-value-list,
    destructuring-bind,
    handler-case, handler-bind, error, warn, signal,
    with-open-file, open, close,
    format, print, princ, prin1, terpri, fresh-line,
    read, read-line, read-char,
    car, cdr, cons, list, append, reverse, length,
    first, second, third, fourth, fifth,
    rest, last, nth, nthcdr, butlast,
    null, atom, listp, consp, numberp, stringp,
    symbolp, functionp, integerp, floatp, characterp,
    eq, eql, equal, equalp,
    quote, backquote,
    make-hash-table, gethash, remhash, maphash, hash-table-count,
    make-array, aref, array-dimensions,
    make-instance, slot-value, initialize-instance,
    in-package, defpackage, use-package, export, import,
    compile, load, require, provide,
    t, nil%
  },
  morekeywords=[2]{%
    +, -, *, /, mod, rem, expt, sqrt, abs,
    floor, ceiling, truncate, round,
    min, max, 1+, 1-,
    =, /=, <, >, <=, >=,
    char, char-code, code-char, char-upcase, char-downcase,
    string, string-upcase, string-downcase, string-length,
    coerce, type-of, typep, check-type%
  }
}

\lstdefinestyle{lispstyle}{%
  language=CommonLisp,
  basicstyle=\ttfamily\small,
  keywordstyle=[1]{\color{clKeyword}\bfseries},
  keywordstyle=[2]{\color{clKeyword}},
  commentstyle=\itshape\color{clComment},
  stringstyle=\color{clString},
  backgroundcolor=\color{clBg},
  frame=single,
  framesep=5pt,
  framerule=0.4pt,
  rulecolor=\color{gray!40},
  xleftmargin=10pt,
  xrightmargin=4pt,
  breaklines=true,
  breakatwhitespace=true,
  showstringspaces=false,
  tabsize=2,
  numbers=none,
  upquote=true,
  extendedchars=true,
  literate=
    {->}{{$\rightarrow$}}2
    {=>}{{$\Rightarrow$}}2
    {à}{{\`a}}1  {è}{{\`e}}1  {é}{{\'{e}}}1  {ì}{{\`{\i}}}1
    {ò}{{\`o}}1  {ù}{{\`u}}1
}

\lstset{style=lispstyle}
\newcommand{\cl}[1]{\lstinline[style=lispstyle]{#1}}
\newcommand{\fn}[1]{\texttt{\color{clKeyword}\bfseries#1}}

% --- tcolorbox ---
\usepackage[most]{tcolorbox}
\tcbuselibrary{listings, breakable, skins}

% ================================================================
%  BOX TIPOLOGIE
% ================================================================

% ---- CONCETTO ----
% \Needspace{5\baselineskip} = salta pagina se mancano meno di 5 righe
\newenvironment{concetto}[1]{%
  \Needspace{6\baselineskip}%
  \medskip
  \begin{tcolorbox}[
    enhanced,
    colframe=boxConFrame, colback=boxConBack,
    fonttitle=\bfseries\normalsize,
    title={\textcolor{white}{\strut #1}},
    coltitle=white,
    attach boxed title to top left={yshift=-2mm, xshift=6pt},
    boxed title style={
      colback=boxConFrame, colframe=boxConFrame,
      arc=2pt, boxrule=0pt,
      left=6pt, right=6pt, top=2pt, bottom=2pt
    },
    arc=4pt, boxrule=1pt,
    left=12pt, right=12pt, top=14pt, bottom=10pt,
    before skip=18pt plus 4pt, after skip=14pt plus 3pt,
    parskip=6pt
  ]
  \small\color{black!80}%
}{\end{tcolorbox}\medskip}

% ---- SINTASSI ----
\newenvironment{sintassi}{%
  \smallskip
  \begin{tcolorbox}[
    enhanced,
    colframe=csBlue!60, colback=csBlue!4,
    arc=2pt, boxrule=0.6pt,
    left=10pt, right=8pt, top=5pt, bottom=5pt,
    before skip=6pt, after skip=8pt
  ]
  \ttfamily\small\color{csBlue!80!black}%
}{\end{tcolorbox}\smallskip}

% ---- ESEMPIO ----
% breakable ma con soglia minima di righe prima dello spacco
\newenvironment{esempio}[1][Esempio]{%
  \medskip
  \begin{tcolorbox}[
    enhanced, breakable,
    lines before break=5,
    colframe=boxEsFrame, colback=boxEsBack,
    fonttitle=\bfseries\small,
    title={#1},
    coltitle=boxEsFrame!80!black,
    attach boxed title to top left={yshift=-2mm, xshift=6pt},
    boxed title style={
      colback=boxEsBack, colframe=boxEsFrame,
      arc=2pt, boxrule=0.6pt,
      left=5pt, right=5pt, top=1pt, bottom=1pt
    },
    arc=3pt, boxrule=0.8pt,
    left=10pt, right=10pt, top=12pt, bottom=8pt,
    before skip=10pt plus 2pt, after skip=10pt plus 2pt,
    parskip=4pt
  ]
  \small\color{black!80}%
}{\end{tcolorbox}\medskip}

% ---- ATTENZIONE ----
\newenvironment{attenzione}{%
  \Needspace{5\baselineskip}%
  \medskip
  \begin{tcolorbox}[
    enhanced,
    colframe=boxWarnFrame, colback=boxWarnBack,
    fonttitle=\bfseries\small,
    title={$\blacktriangleright$\enspace Attenzione},
    coltitle=boxWarnFrame!80!black,
    attach boxed title to top left={yshift=-2mm, xshift=6pt},
    boxed title style={
      colback=boxWarnBack, colframe=boxWarnFrame,
      arc=2pt, boxrule=0.6pt,
      left=5pt, right=5pt, top=1pt, bottom=1pt
    },
    arc=3pt, boxrule=0.8pt,
    left=10pt, right=10pt, top=12pt, bottom=8pt,
    before skip=10pt plus 2pt, after skip=10pt plus 2pt,
    parskip=4pt
  ]
  \small\color{black!80}%
}{\end{tcolorbox}\medskip}

% ---- INTUIZIONE ----
\newenvironment{intuizione}{%
  \Needspace{5\baselineskip}%
  \medskip
  \begin{tcolorbox}[
    enhanced,
    colframe=boxKeyFrame, colback=boxKeyBack,
    fonttitle=\bfseries\small,
    title={$\bigstar$\enspace Intuizione chiave},
    coltitle=boxKeyFrame!80!black,
    attach boxed title to top left={yshift=-2mm, xshift=6pt},
    boxed title style={
      colback=boxKeyBack, colframe=boxKeyFrame,
      arc=2pt, boxrule=0.6pt,
      left=5pt, right=5pt, top=1pt, bottom=1pt
    },
    arc=3pt, boxrule=0.8pt,
    left=12pt, right=12pt, top=12pt, bottom=10pt,
    before skip=10pt plus 2pt, after skip=10pt plus 2pt,
    parskip=4pt
  ]
  \small\color{black!80}%
}{\end{tcolorbox}\medskip}

% ---- WORKBOOK ----
\newenvironment{workbook}{%
  \Needspace{6\baselineskip}%
  \medskip
  \begin{tcolorbox}[
    enhanced,
    colframe=boxWbFrame, colback=boxWbBack,
    fonttitle=\bfseries\small,
    title={$\blacktriangleright$\enspace Nel workbook},
    coltitle=boxWbFrame!80!black,
    attach boxed title to top left={yshift=-2mm, xshift=6pt},
    boxed title style={
      colback=boxWbBack, colframe=boxWbFrame,
      arc=2pt, boxrule=0.6pt,
      left=5pt, right=5pt, top=1pt, bottom=1pt
    },
    arc=3pt, boxrule=0.8pt,
    left=12pt, right=12pt, top=12pt, bottom=8pt,
    before skip=10pt plus 2pt, after skip=10pt plus 2pt,
    parskip=3pt
  ]
  \small\color{black!80}%
}{\end{tcolorbox}\medskip}

% ---- NOTA ----
\newenvironment{nota}{%
  \Needspace{4\baselineskip}%
  \smallskip
  \begin{tcolorbox}[
    enhanced,
    colframe=boxNoteFrame, colback=boxNoteBack,
    arc=3pt, boxrule=0.5pt,
    left=10pt, right=10pt, top=6pt, bottom=6pt,
    before skip=8pt, after skip=8pt
  ]
  \small\color{black!65}\itshape%
}{\end{tcolorbox}\smallskip}

% ---- CONFRONTO ----
\definecolor{boxCfFrame}{RGB}{160,80,20}
\definecolor{boxCfBack}{RGB}{253,247,238}
\newenvironment{confronto}{%
  \Needspace{4\baselineskip}%
  \smallskip
  \begin{tcolorbox}[
    enhanced,
    colframe=boxCfFrame, colback=boxCfBack,
    leftrule=3pt, toprule=0.3pt, bottomrule=0.3pt, rightrule=0.3pt,
    arc=2pt,
    left=12pt, right=10pt, top=6pt, bottom=6pt,
    before skip=8pt, after skip=8pt
  ]
  \small\color{black!70}%
}{\end{tcolorbox}\smallskip}

% ================================================================
%  SEPARATORI E COMANDI SEMANTICI
% ================================================================

% Separatore orizzontale leggero
\newcommand{\seprule}{%
  \par\vspace{6pt}%
  \noindent\textcolor{black!20}{\rule{\linewidth}{0.4pt}}%
  \par\vspace{4pt}%
}

% Titolo di argomento interno (sotto-sezione visiva, senza numero)
\newcommand{\argomento}[1]{%
  \Needspace{5\baselineskip}%
  \par\vspace{16pt}%
  \noindent{\large\bfseries\color{csBlue} #1}%
  \par\noindent\textcolor{csBlue!35}{\rule{\linewidth}{0.7pt}}%
  \par\vspace{4pt}%
}

% Testo descrittivo/contestuale: colore leggermente attenuato, size normale
% Usato per paragrafi di contesto non essenziali
\newcommand{\testo}[1]{%
  \par\smallskip%
  {\small\color{black!65} #1}%
  \par\smallskip%
}

% Label inline per concetti chiave dentro il testo
\newcommand{\concettoinline}[1]{%
  \textbf{\textcolor{csBlue}{#1}}%
}

% ---- Comandi per gerarchia INTERNA ai box ----

% Frase-chiave: grande, colorata, su riga propria
\newcommand{\keyline}[1]{%
  \par\smallskip%
  {\normalsize\bfseries\color{csBlue!90!black} #1}%
  \par\vspace{3pt}%
}

% Dettaglio: testo attenuato, leggermente più piccolo
\newcommand{\detail}[1]{%
  {\small\color{black!60} #1}%
}

% Separatore interno leggero (dentro i box)
\newcommand{\boxsep}{%
  \par\vspace{4pt}%
  \noindent\textcolor{black!15}{\rule{\linewidth}{0.3pt}}%
  \par\vspace{4pt}%
}

% Voce di elenco con etichetta colorata inline
% \boxitem{Etichetta}{testo descrittivo}
\newcommand{\boxitem}[2]{%
  \par\smallskip%
  \noindent{\small\bfseries\color{csBlue!80!black}#1}\enspace%
  {\small\color{black!70}#2}%
}

% Comando per funzioni inline
\newcommand{\repl}[1]{\texttt{\small\color{csGreen!70!black}#1}}
\newcommand{\risultato}[1]{\texttt{\small\color{csGreen!70!black}$\Rightarrow$\enspace #1}}

% Nota inline a piè di paragrafo (sostituisce marginnote)
\newcommand{\mnota}[1]{%
  \par\smallskip\noindent%
  \textcolor{black!45}{\footnotesize\itshape\hspace{1em}$\triangleright$\enspace #1}%
  \par\smallskip%
}

% ================================================================
%  INTESTAZIONI DI PAGINA
% ================================================================

\usepackage{fancyhdr}
\pagestyle{fancy}
\fancyhf{}
% Solo il titolo corrente del capitolo a sinistra, pagina a destra
% Troncato per evitare overflow
\fancyhead[L]{\small\color{black!50}\itshape\nouppercase{\leftmark}}
\fancyhead[R]{\small\color{black!40}\thepage}
\fancyfoot{}
\renewcommand{\headrulewidth}{0.3pt}
\renewcommand{\footrulewidth}{0pt}

% Ridefinisce \chaptermark per non stampare "CAPITOLO N" ma solo il titolo
\renewcommand{\chaptermark}[1]{\markboth{\chaptername~\thechapter.\ #1}{}}

% ================================================================
%  STILE CAPITOLI E SEZIONI
% ================================================================

\usepackage{titlesec}

% Capitoli: numero grande colorato + titolo + riga
\titleformat{\chapter}[display]
  {\normalfont\huge\bfseries\color{csBlue}}
  {\normalfont\large\color{csBlue!55}\scshape\chaptername\ \thechapter}
  {6pt}
  {\Huge\color{csBlue}}
  [\vspace{4pt}{\color{csBlue!25}\rule{\linewidth}{1.5pt}}\vspace{2pt}]

\titlespacing*{\chapter}{0pt}{0pt}{28pt}

% Sezioni: numero colorato + titolo + sottile riga
\titleformat{\section}
  {\normalfont\large\bfseries\color{black!85}}
  {\textcolor{csBlue}{\thesection.}}
  {7pt}
  {}
  [\vspace{1pt}{\color{csBlue!25}\rule{\linewidth}{0.5pt}}]

\titlespacing*{\section}{0pt}{22pt}{8pt}

% Sottosezioni
\titleformat{\subsection}
  {\normalfont\normalsize\bfseries\color{csBlue!80!black}}
  {\thesubsection.}
  {6pt}
  {}

\titlespacing*{\subsection}{0pt}{14pt}{4pt}

% ================================================================
%  HYPERREF (sempre ultimo)
% ================================================================

\usepackage[
  colorlinks=true,
  linkcolor=csBlue,
  urlcolor=csTeal,
  citecolor=csGreen,
  pdftitle={Common Lisp - Fondamenti e Pratica},
  pdfauthor={Carmelo Calzerano}
]{hyperref}

% ================================================================
%  PACCHETTI ACCESSORI
% ================================================================

\usepackage{enumitem}
\setlist[itemize]{leftmargin=1.6em, itemsep=2pt, topsep=3pt}
\setlist[enumerate]{leftmargin=1.6em, itemsep=2pt, topsep=3pt}

\usepackage{booktabs}
\usepackage{array}
\usepackage{multicol}
\usepackage{ragged2e}
\usepackage{needspace}  % per evitare pagebreak interni ai box brevi
